% GENERATED FILE. Edit canonical lessons, then run npm run generate:books.
% canonical-source-hash: fnv1a64:73982ae907ded5a0
% canonical-lessons: FA-C99-polow, FA-C99-ash, FA-C99-hendevane, FA-C99-kharboze, FA-C99-holu

\chapter{Rice pilaf, Thick soup, Watermelon, Melon, Peach}
\label{ch:fa-rice-pilaf-thick-soup-watermelon-melon-peach}
\hlchaptermodality{99}

\begin{chapteropening}
\textbf{By the end of this chapter:} \emph{I can say rice pilaf, a thick soup, a watermelon, a melon and a peach.}

Complete the last lesson of chapter 99: I can say rice pilaf, a thick soup, a watermelon, a melon and a peach.
\end{chapteropening}

\section[\texorpdfstring{\fa{پلو} (polow)}{پلو (polow)}]{\fa{پلو} (polow) — rice pilaf}

\label{lesson:FA-C99-polow}



\begin{quote}
Before the new one: say the Persian for a salad, then the Persian for a kebab.
\end{quote}

\subsection*{You'll want to know: \fa{پلو}}
\textbf{\fa{پلو}} — \emph{polow} — \textquotedblleft{}rice pilaf\textquotedblright{}.

\begin{cousinweb}
English \emph{pilaf} comes from this word through Turkish; the origin of \emph{polow} itself is not settled.
\end{cousinweb}

\begin{grammarlens}[title={What you've built}]
One more word for this chapter.
\end{grammarlens}

\subsection*{Your turn}
\begin{itemize}
\raggedright
  \item \emph{Say it:} \emph{polow}
  \item \emph{Say it:} \emph{polow}, once more
  \item \emph{Say it:} say \emph{polow}
  \item \emph{Recall:} say the Persian for a salad, then the Persian for a kebab, then say \emph{polow} again
\end{itemize}

\begin{tcolorbox}[breakable,colback=teal!4,colframe=teal!35!black,title={Before you move on}]
What does \fa{پلو} mean? (\textquotedblleft{}Rice pilaf\textquotedblright{}.) Say it once more.
\end{tcolorbox}
\section[\texorpdfstring{\fa{آش} (âsh)}{آش (âsh)}]{\fa{آش} (âsh) — a thick soup}

\label{lesson:FA-C99-ash}



\begin{quote}
Before the new one: say the Persian for a kebab, then the Persian for rice pilaf.
\end{quote}

\subsection*{You'll want to know: \fa{آش}}
\textbf{\fa{آش}} — \emph{âsh} — \textquotedblleft{}a thick soup\textquotedblright{}.

\begin{cousinweb}
\emph{âsh} is sometimes linked with Turkic \emph{aš}, \textquotedblleft{}food\textquotedblright{}, but no source used here confirms it, and its origin is not settled.
\end{cousinweb}

\begin{grammarlens}[title={What you've built}]
One more word for this chapter.
\end{grammarlens}

\subsection*{Your turn}
\begin{itemize}
\raggedright
  \item \emph{Say it:} \emph{âsh}
  \item \emph{Say it:} \emph{âsh}, once more
  \item \emph{Say it:} say \emph{âsh}
  \item \emph{Recall:} say the Persian for a kebab, then the Persian for rice pilaf, then say \emph{âsh} again
\end{itemize}

\begin{tcolorbox}[breakable,colback=teal!4,colframe=teal!35!black,title={Before you move on}]
What does \fa{آش} mean? (\textquotedblleft{}A thick soup\textquotedblright{}.) Say it once more.
\end{tcolorbox}
\section[\texorpdfstring{\fa{هندوانه} (hendevâne)}{هندوانه (hendevâne)}]{\fa{هندوانه} (hendevâne) — a watermelon}

\label{lesson:FA-C99-hendevane}



\begin{quote}
Before the new one: say the Persian for rice pilaf, then the Persian for a thick soup.
\end{quote}

\subsection*{You'll want to know: \fa{هندوانه}}
\textbf{\fa{هندوانه}} — \emph{hendevâne} — \textquotedblleft{}a watermelon\textquotedblright{}.

\begin{cousinweb}
\emph{hendevâne} is usually explained as \textquotedblleft{}the Indian one\textquotedblright{}, built on \emph{Hendu}, \textquotedblleft{}India, Indian\textquotedblright{}, with the ending \emph{-âne}: the watermelon named for the land it was thought to come from.
\end{cousinweb}

\begin{grammarlens}[title={What you've built}]
One more word for this chapter.
\end{grammarlens}

\subsection*{Your turn}
\begin{itemize}
\raggedright
  \item \emph{Say it:} \emph{hendevâne}
  \item \emph{Say it:} \emph{hendevâne}, once more
  \item \emph{Say it:} say \emph{hendevâne}
  \item \emph{Recall:} say the Persian for rice pilaf, then the Persian for a thick soup, then say \emph{hendevâne} again
\end{itemize}

\begin{tcolorbox}[breakable,colback=teal!4,colframe=teal!35!black,title={Before you move on}]
What does \fa{هندوانه} mean? (\textquotedblleft{}A watermelon\textquotedblright{}.) Say it once more.
\end{tcolorbox}
\section[\texorpdfstring{\fa{خربزه} (kharboze)}{خربزه (kharboze)}]{\fa{خربزه} (kharboze) — a melon}

\label{lesson:FA-C99-kharboze}



\begin{quote}
Before the new one: say the Persian for a thick soup, then the Persian for a watermelon.
\end{quote}

\subsection*{You'll want to know: \fa{خربزه}}
\textbf{\fa{خربزه}} — \emph{kharboze} — \textquotedblleft{}a melon\textquotedblright{}.

\begin{cousinweb}
\emph{kharboze} is usually traced to Middle Persian \emph{xarbūzag}, \textquotedblleft{}melon\textquotedblright{}; how that word was built is not settled.
\end{cousinweb}

\begin{grammarlens}[title={What you've built}]
One more word for this chapter.
\end{grammarlens}

\subsection*{Your turn}
\begin{itemize}
\raggedright
  \item \emph{Say it:} \emph{kharboze}
  \item \emph{Say it:} \emph{kharboze}, once more
  \item \emph{Say it:} say \emph{kharboze}
  \item \emph{Recall:} say the Persian for a thick soup, then the Persian for a watermelon, then say \emph{kharboze} again
\end{itemize}

\begin{tcolorbox}[breakable,colback=teal!4,colframe=teal!35!black,title={Before you move on}]
What does \fa{خربزه} mean? (\textquotedblleft{}A melon\textquotedblright{}.) Say it once more.
\end{tcolorbox}
\section[\texorpdfstring{\fa{هلو} (holu)}{هلو (holu)}]{\fa{هلو} (holu) — a peach}

\label{lesson:FA-C99-holu}



\begin{quote}
Before the new one: say the Persian for a watermelon, then the Persian for a melon.
\end{quote}

\subsection*{You'll want to know: \fa{هلو}}
\textbf{\fa{هلو}} — \emph{holu} — \textquotedblleft{}a peach\textquotedblright{}.

\begin{cousinweb}
\emph{holu} is the everyday word for a peach, but no source used here gives its origin, which is not settled.
\end{cousinweb}

\begin{grammarlens}[title={What you've built}]
One more word for this chapter.
\end{grammarlens}

\subsection*{Your turn}
\begin{itemize}
\raggedright
  \item \emph{Say it:} \emph{holu}
  \item \emph{Say it:} \emph{holu}, once more
  \item \emph{Say it:} say \emph{holu}
  \item \emph{Recall:} say the Persian for a watermelon, then the Persian for a melon, then say \emph{holu} again
\end{itemize}

\begin{tcolorbox}[breakable,colback=teal!4,colframe=teal!35!black,title={Before you move on}]
What does \fa{هلو} mean? (\textquotedblleft{}A peach\textquotedblright{}.) Say it once more.
\end{tcolorbox}
